\chapter{Conclusion}
\label{ch:conclusion}

This thesis measured static navigation extraction on six Angular applications, classified
its losses and evaluated scope attribution. Table~\ref{tab:rqanswers} summarises the answers
to the five research questions.

\begingroup\footnotesize
\begin{xltabular}{\textwidth}{@{}>{\raggedright\arraybackslash}p{0.13\textwidth}YYY@{}}
\caption{Answers to the research questions, supporting figures and main limitations.}
\label{tab:rqanswers}\\
\toprule
RQ & Answer & Key figures & Main limit \\
\midrule
\endfirsthead
\toprule
RQ & Answer & Key figures & Main limit \\
\midrule
\endhead
\bottomrule
\endfoot
RQ1: how much is recovered
& Every declared screen identified by the static readers has a node; few constructs yield must transitions.
& Nodes complete against the mounted routes on \ProbeExactApps{} of six applications. \CensusMustSites{} constructs yield \CensusMust{} must edges, \CensusMustKdash{} from one toolbar; \RouteOrderMustWithout{} without an assumed route order.
& Most must edges depend on the unvalidated route-order assumption (\RouteOrderResting{} edges); the toolbar's scope is untested. \\
RQ2: where extraction fails
& Five defect classes and one candidate class (Table~\ref{tab:defects}); attribution is the largest loss category in every application under the chosen rule order.
& Table~\ref{tab:loss}; testing the target first makes the target category larger than attribution on \AttrOrderFlipWhere{}.
& The ranking depends on the order in which the loss rules are applied. \\
RQ3: what scope attribution adds
& It recovers constructs without losing any recovered by the baseline, mostly through earlier mechanisms and almost all as possible transitions. On the alfresco reference, the added step trades precision for recall.
& \AttrRecoveredShare{} of the attribution constructs (\AttrRecoveredShareNoMifos{} without mifos), with screen sets \BaselineFactorMin{} to \BaselineFactorMax{} times narrower than every represented site on every screen. The step adds \AblationFullAddedSites{} sites (\AblationFullAddedShare{}) over \emph{outlet}. Alfresco reference: recall \GTAlfFullRecall{} against \GTAlfClassRecall{} for a per-class join, precision \GTAlfFullPrecision{} against \GTAlfClassPrecision{}.
& Both structural references apply the analysis's own relations. The added step recovers no extra end-to-end transition over \emph{outlet} on either application with a suite (mifos: \EtoEMifosHitOutlet{}/\EtoEMifosPairs{} with and without it, against \EtoEMifosChanceHits{} for uniformly drawn pairs at the graph's density). The comparison with Rima et al.'s tool gives no overall ranking (Section~\ref{sec:res-h2h}). \\
RQ4: whether the graph responds to change
& A conditional sanity check of witness mapping on three applications; change-impact accuracy remains unmeasured.
& \MutDetected{}/\MutScored{} scored mutants detected.
& Operators act only at constructs already represented, making detection close to guaranteed; the harness was corrected after inspecting its misses. \\
RQ5: whose losses they are
& Three groups emerge: repaired tool defects; unresolved cases (remaining attribution losses, named outlets and one application's computed routes); and cases beyond static route-table analysis (a custom \texttt{UrlMatcher} and deployment-selected routes). Application decisions trigger may classification under the present demotion rules; this does not establish an inherent limit of static analysis.
& \MayAppAny{}/\MayOnlyConstructs{} may-only witness locations sit behind an application decision, \MayToolOnly{} behind the extractor's uncertainty alone; \AttrRemainingUnexplained{} of the \AttrRemaining{} remaining attribution constructs are labelled with no reason found.
& These causes describe the classes found and repaired. The harness enforces the presence of a label, not its correctness. \\
\end{xltabular}
\endgroup

\paragraph{Contributions.} The primary contribution is the measurement: six pinned applications,
five evaluation methods with different blind spots, four further scope checks and a classification
of every lost construct. The extractor supports nested and parameterised route trees, which the
stated restrictions of an earlier Angular extraction exclude. Those restrictions admit
\EarlierWorkAdmittedApps{} of the \EarlierWorkApps{} corpus applications; this extractor processes
all of them (Section~\ref{sec:res-earlier}).

The evaluation apparatus makes the measurements inspectable. Tables are generated from harness
results, and source-based methods re-run from pinned commits. Runtime reproduction is partial:
the route oracle re-runs only in part, the realworld end-to-end check uses an unpinned public
backend, and superseded figures cannot be regenerated.

Scope attribution is specified as a dataflow analysis over screen sets and implemented over
three of its five input relations (Section~\ref{sec:attribution}). Most recovery comes from earlier
mechanisms included as analysis levels. The added step places navigation in services, guards
and effects. Its type-directed, transitive service join generalises the bounded, one-level,
name-matched router-service tracing in Rima et al.'s implementation. With a fallback join and
stand-ins for guards and effects, it adds \AblationFullAddedSites{} sites
(\AblationFullAddedShare{}) over outlet descent. The call graph alone adds
\AblationCallGraphNetSites{} (\AblationCallGraphNetShare{}) over the per-class rule it replaced
(Section~\ref{sec:res-ablation}).

\paragraph{What remains open.} Dispatch is specified but unimplemented, guards use a stand-in,
and the call graph ignores the injector hierarchy. Screen sets were checked in both directions
on only two applications, against references that apply the analysis's own relations.

The main corpus is in-sample. On Rima et al.'s seven applications, excluded from initial
development, the extractor contains \RimaRepHitFull{} of the \RimaPairs{} reference pairs after
four repairs (\RimaHitFull{} before them). These results are therefore partly fitted too
(Section~\ref{sec:res-rima}). The comparison with their tool uses small AI-read samples; I
confirmed \HHAuthorAgree{} of the \HHAuthorJudged{} cases I re-read
(Section~\ref{sec:res-h2h}). The mifos end-to-end check exposed unsupported target forms and
a relative link resolved against a child tab's address instead of its containing route.
Repairs after that check make the final mifos result in-sample (Section~\ref{sec:res-e2e}).

The must-tier judgments were made by the AI assistant; two readings by one model measure
consistency rather than correctness. Rendered-interface checks covered only screens reachable
without signing in. Signed-in execution is needed to test the rest (Section~\ref{sec:future}).
Consumer benefit, including support for UI agents and test generation, remains unmeasured.

A route table alone is insufficient for static navigation extraction. May edges are candidates
that require further checks. Must transitions (\CensusMustSites{} constructs) hold only under
the assumptions of Claim~\ref{claim:must}.
